
\documentclass{article}
\usepackage{colortbl}
\usepackage{makecell}
\usepackage{multirow}
\usepackage{supertabular}

\begin{document}

\newcounter{utterance}

\centering \large Interaction Transcript for game `bbh', experiment `hyperbaton', episode 55 with Bentel rockers\_\_qwen3.5{-}4b\_\_Bentel\_iter\_3.
\vspace{24pt}

{ \footnotesize  \setcounter{utterance}{1}
\setlength{\tabcolsep}{0pt}
\begin{supertabular}{c@{$\;$}|p{.15\linewidth}@{}p{.15\linewidth}p{.15\linewidth}p{.15\linewidth}p{.15\linewidth}p{.15\linewidth}}
    \# & \multicolumn{2}{c}{Player} && \multicolumn{2}{c}{Game Master} \\
    \hline

    \theutterance \stepcounter{utterance}  
    & & & \multicolumn{4}{p{0.6\linewidth}}{
        \cellcolor[rgb]{0.9,0.9,0.9}{
            \makecell[{{p{\linewidth}}}]{
                \texttt{\tiny{[P1$\langle$GM]}}
                \texttt{Order adjectives correctly in English sentences.} \\
\\ 
\texttt{Q: Which sentence has the correct adjective order:} \\
\texttt{Options:} \\
\texttt{(A) rubber terrible ship} \\
\texttt{(B) terrible rubber ship} \\
\texttt{A: (B)} \\
\\ 
\texttt{Q: Which sentence has the correct adjective order:} \\
\texttt{Options:} \\
\texttt{(A) repulsive small Brazilian exercise ship} \\
\texttt{(B) Brazilian repulsive exercise small ship} \\
\texttt{A: (A)} \\
\\ 
\texttt{Q: Which sentence has the correct adjective order:} \\
\texttt{Options:} \\
\texttt{(A) blue gold wonderful square shoe} \\
\texttt{(B) wonderful square blue gold shoe} \\
\texttt{A: (B)} \\
\\ 
\texttt{Q: Which sentence has the correct adjective order:} \\
\texttt{Options:} \\
\texttt{(A) silly huge old{-}fashioned triangular tan Pakistani baby} \\
\texttt{(B) triangular tan Pakistani silly huge old{-}fashioned baby} \\
\texttt{A:} \\
            }
        }
    }
     \\ \\

    \theutterance \stepcounter{utterance}  
    & \multicolumn{4}{p{0.6\linewidth}}{
        \cellcolor[rgb]{0.561,0.737,0.561}{
            \makecell[{{p{\linewidth}}}]{
                \texttt{\tiny{[P1$\rangle$GM]}}
                \texttt{(A)} \\
\\ 
\texttt{**Explanation:**} \\
\texttt{The correct order of adjectives in English generally follows this pattern:} \\
\texttt{1.  **Opinion** (e.g., silly)} \\
\texttt{2.  **Size** (e.g., huge)} \\
\texttt{3.  **Age** (e.g., old{-}fashioned)} \\
\texttt{4.  **Shape** (e.g., triangular)} \\
\texttt{5.  **Color** (e.g., tan)} \\
\texttt{6.  **Origin** (e.g., Pakistani)} \\
\texttt{7.  **Material** (none in this specific phrase, but "baby" acts as the noun here, though "baby" is actually a noun, the other adjectives modify it)} \\
\\ 
\texttt{Let's re{-}evaluate the specific phrase "silly huge old{-}fashioned triangular tan Pakistani baby":} \\
\texttt{{-}   **Opinion:** silly} \\
\texttt{{-}   **Size:** huge} \\
\texttt{{-}   **Age:** old{-}fashioned} \\
\texttt{{-}   **Shape:** triangular} \\
\texttt{{-}   **Color:** tan} \\
\texttt{{-}   **Origin:** Pakistani} \\
\texttt{{-}   **Noun:** baby} \\
\\ 
\texttt{This follows the standard sequence perfectly. Option (B) scrambles these elements, placing the origin before the opinion and shape, which is incorrect.} \\
\\ 
\texttt{Therefore, the correct choice is **(A)**.} \\
            }
        }
    }
    & & \\ \\

    \theutterance \stepcounter{utterance}  
    & & & \multicolumn{2}{p{0.3\linewidth}}{
        \cellcolor[rgb]{0.9,0.9,0.9}{
            \makecell[{{p{\linewidth}}}]{
                \texttt{\tiny{[GM$|$GM]}}
                \texttt{Target: (A)} \\
            }
        }
    }
    & & \\ \\

    \theutterance \stepcounter{utterance}  
    & & & \multicolumn{2}{p{0.3\linewidth}}{
        \cellcolor[rgb]{0.9,0.9,0.9}{
            \makecell[{{p{\linewidth}}}]{
                \texttt{\tiny{[GM$|$GM]}}
                \texttt{game\_result = LOSE} \\
            }
        }
    }
    & & \\ \\

\end{supertabular}
}

\end{document}
